\section{Introduction}\label{sec:intro}

French residential real estate is an idiosyncratic asset class for
retail investors. The price process mean-reverts slowly; the rental
cash flow is heterogeneous and exposed to tenant-specific default risk;
and the legal framework around eviction, the \textit{Indice de Référence
des Loyers} and the Visale guarantee shapes realised loss in ways that
do not map cleanly onto traded credit instruments.

This paper asks whether, given the same underlying building, a tranched
cash-flow structure can deliver a different risk-return profile from
outright ownership or a classical SAS pool. We formalise three
investment vehicles --- Model A: single-apartment ownership; Model B: a
pro-rata SAS share; Model C: a CDO-style senior / mezzanine / equity
tranche stack --- and compute their fair price and risk metrics on the
same set of Monte Carlo paths.

The simulation pipeline is calibrated on \dataParisN{} quarters of
Notaires-INSEE Paris price data and \dataOatN{} months of OAT 10Y data;
every numerical estimate quoted in the paper is traceable to a row of
\texttt{artifacts/results.csv} produced by the same code. The credit
layer juxtaposes the Gaussian one-factor copula
\autocite{Li2000, Vasicek2002}, the Student-t copula
\autocite{EmbrechtsMcNeilStraumann2002}, and a Cox doubly stochastic
intensity \autocite{DuffieSingleton2003}, so the tail dependence that
the Gaussian benchmark silently sets to zero becomes visible at the
senior attach point. The rental-default insurance is priced both by the
actuarial standard-deviation principle and as a put on rent shortfall,
and the gap between the two is reported.

\subsection{Literature context}

The CDO tranching mechanic borrows from two modelling traditions. The
structural side originates with \autocite{BlackScholes1973} and
\autocite{Merton1976}: a firm defaults when its asset value crosses a
threshold, with the asset following a geometric Brownian motion. The
reduced-form side, started by \autocite{DuffieSingleton2003} and
expanded by \autocite{Schonbucher2003}, treats default as the first
jump of a Cox process driven by an exogenous hazard. The Cox
doubly-stochastic credit model in \cref{ssec:cox} is the canonical
example of the reduced-form mechanic with a CIR-like macro factor.

Copula-based pricing of multi-name credit baskets took off with
\autocite{Li2000}, which mapped default times into a joint Gaussian
distribution; \autocite{Vasicek2002} provided the large-portfolio
closed form that became the industry workhorse for CDO tranche pricing.
\autocite{EmbrechtsMcNeilStraumann2002} pointed out that the Gaussian
copula misses tail dependence, and \autocite{DonnellyEmbrechts2010}
quantified the implication for senior-tranche pricing in the subprime
crisis. \autocite{MacKenzieSpears2014} traces the Li (2000) formula's
twin role as industry standard and post-crisis scapegoat. The
Student-t copula benchmark in \cref{ssec:t-copula} is the alternative
that this strand of literature recommends.

The waterfall mechanic follows \autocite{AndersenSideniusBasu2003}: a
notional-depleting senior / mezzanine / equity stack with separate
interest and principal cascades. The contribution of the paper is not
in the modelling primitives, which are standard; it is in the
consistent calibration of those primitives on the public French
residential-rent dataset, juxtaposed with two retail benchmarks
(single-apartment ownership and a SAS pooled share) on a shared Monte
Carlo seed. The variance-reduction tooling follows
\autocite{Glasserman2003} for antithetic sampling and
\autocite{Owen1997} for scrambled Sobol QMC; risk measurement uses the
coherent expected-shortfall framework of \autocite{AcerbiTasche2002}.
